% =========================================================================================
% Chapter 6 -- Clinical Reasoning Agent
%
% Numbers here were read from the code and from models/clinical_reasoning_rag_ab/.
% Regenerate the A/B figures with notebooks/clinical_reasoning_rag_ab.ipynb.
% =========================================================================================
\section{Clinical Reasoning Agent}
\label{sec:reasoning}

\subsection{Introduction}

The perception modules described in Chapter~\ref{sec:perception} answer questions about
images. A clinician asks a different question: what is wrong with this patient, and what
should be done next. No combination of image classifiers answers that, because the answer
also depends on the vitals, the laboratory results and the history, none of which the
imaging models can see.

This chapter describes the layer that brings those pieces together. It has one organising
idea, stated below, and one design rule: the decisions that matter for safety are computed
by ordinary code, not by the language model.

\begin{keystatement}
Absent is not normal. A test that was never performed cannot support or contradict any
diagnosis, and a model with no healthy class can never exclude disease.
\end{keystatement}

\subsection{The Clinical State}
\label{subsec:clinical-state}

Everything the agent knows about a patient is collected into a single structured record
called the clinical state. It holds three kinds of information: the imaging findings, the
measured values, and --- importantly --- what is missing.

Imaging findings are recorded in three states rather than two. A finding may be
\emph{detected}, meaning the module looked and saw it; \emph{not detected}, meaning the
module looked and did not see it; or simply absent from the record, meaning nobody looked.
The middle case is a real observation and the last one is not. Collapsing them would throw
away information in one direction and invent it in the other.

Vitals and laboratory values are stored with their reference ranges, and each is flagged
high, low or normal by the code before the model sees it. A value that was not measured is
listed in a separate block, printed under the heading \texttt{NOT MEASURED --- absent, NOT
normal}. Tab.~\ref{tab:statefields} summarises the record.

\begin{table}[H]
  \centering
  \small
  \begin{tabular}{lp{9.2cm}}
    \toprule
    \textbf{Part of the state} & \textbf{What it holds} \\
    \midrule
    Imaging findings   & label, confidence, detected or not detected, evidence grade \\
    Organs not assessed& organs that were requested but never scanned \\
    Vitals             & value, reference range, high/low/normal flag \\
    Laboratory values  & value, reference range, high/low/normal flag \\
    Missing            & every vital and laboratory value that was not measured \\
    Conflicts          & disagreements between the triage and imaging agents \\
    Case quality       & an overall grade with the reasons that produced it \\
    \bottomrule
  \end{tabular}
  \caption{The clinical state passed to the reasoning layer.}
  \label{tab:statefields}
\end{table}

The two perception agents produce differently shaped outputs --- the triage agent gives one
urgency tier with a confidence, the ultrasound agent gives a list of findings --- and both
are normalised into this single form. Each finding also receives an evidence grade:
\emph{experimental} when the module's confidence is not calibrated, \emph{limited} when the
confidence is low or the training evidence was thin, and \emph{strong} otherwise. The grade
travels with the finding so that later decisions can take account of how much the number is
worth.

\subsection{Conflict Detection}
\label{subsec:reasoning-conflict}

The two agents can disagree. The clearest example is a patient triaged as low urgency whose
scan shows a severe finding with high confidence. The system flags this and does not try to
settle it. It has no information that would say which source is wrong, so averaging the two
or preferring the more confident one would be a guess presented as a judgement. The
disagreement is reported and the case is escalated.

\subsection{Escalation}
\label{subsec:reasoning-escalation}

Whether a case is answered directly or escalated is decided before the language model runs,
by the rules in Tab.~\ref{tab:escalation}. Each rule reads only the structured state. None
of them depends on the model having noticed anything, so the escalation behaviour is the
same regardless of which model is used, or whether the model fails entirely.

\begin{table}[H]
  \centering
  \small
  \begin{tabular}{p{10.5cm}}
    \toprule
    \textbf{The case is escalated when\ldots} \\
    \midrule
    the two agents disagree \\
    the overall case quality is graded poor \\
    a high-risk finding rests on evidence that is not strong \\
    a high-risk finding is below the confidence considered decisive (0.80) \\
    imaging is positive but a key laboratory value (troponin, lactate) was not measured \\
    nothing was found, but the models cannot rule out disease \\
    an organ was requested and never assessed \\
    \bottomrule
  \end{tabular}
  \caption{Escalation rules. All are evaluated on the structured state before generation.}
  \label{tab:escalation}
\end{table}

The last rule deserves a note. It applies whether or not something was found elsewhere. A
finding in the lung says nothing about a heart nobody scanned, and a positive result in one
organ raises the importance of the gap rather than closing it.

\subsection{Knowledge Corpus}
\label{subsec:reasoning-kb}

To let the agent consult reference material rather than rely only on what the model
remembers, a small corpus of thirty knowledge units was assembled. Each unit answers one
clinical question and carries the citation it came from. Tab.~\ref{tab:corpus} shows the
composition.

\begin{table}[H]
  \centering
  \small
  \begin{tabular}{llr}
    \toprule
    \textbf{Group} & \textbf{Covers} & \textbf{Units} \\
    \midrule
    Lung       & B-lines, A-lines, lung sliding, pneumothorax, BLUE protocol & 10 \\
    Cardiac    & left and right ventricle, pericardial effusion, tamponade   & 5 \\
    FAST       & the four trauma views and their interpretation              & 5 \\
    Shock      & classification and the four shock categories                & 5 \\
    Boundaries & what a focused study does not establish                     & 5 \\
    \bottomrule
  \end{tabular}
  \caption{The thirty knowledge units, by group.}
  \label{tab:corpus}
\end{table}

All thirty are drawn from six open-access papers under Creative Commons licences, so the
text can be redistributed with the project. Two constraints shaped what went in. First,
laboratory reference ranges were deliberately kept out: comparing a value against a
threshold is a fixed calculation the code already performs, and putting the same
information in a retrieval corpus would make a settled comparison depend on a similarity
score. Second, passages were chosen for the limits they state rather than for definitions.
The most useful units are those recording what a test cannot settle --- that the lung point
sign is only 66\% sensitive, that lung sliding is absent in eight conditions besides
pneumothorax, or that a negative FAST detects free fluid reliably only above
150--200\,mL.

The corpus does not cover obstetric, vascular or paediatric ultrasound, and it is not a
substitute for a guideline. It is small by design: enough to test whether retrieval helps,
not enough to claim coverage of emergency ultrasound.

\subsection{Retrieval}
\label{subsec:reasoning-rag}

Retrieval uses TF--IDF over the thirty units. For a corpus of this size it performs
comparably to a neural retriever, needs no model download, and its scores can be inspected,
which matters more here than a small gain in ranking.

Two rules govern how results are used. A passage below a relevance threshold is reported as
no result at all, because handing the model plausible text that does not bear on the case
creates a new opportunity for unsupported claims. And any citation the model produces must
correspond to a passage that was actually retrieved. If nothing was retrieved, the prompt
says so explicitly rather than staying silent, so the model cannot imply guideline support
it never had.

\subsection{The Language Model}
\label{subsec:reasoning-llm}

The model is HuatuoGPT-o1-8B, a medical model, run locally through \texttt{llama.cpp} in
4-bit quantised form (4.9\,GB). Running locally rather than through an API is a deliberate
choice: patient data does not leave the machine.

The cost is real. On a consumer CPU a single case takes roughly 23 minutes. On a Colab T4
the same case takes 40 to 90 seconds, and the full five-case benchmark runs in about five
minutes. Temperature is set to zero with a fixed seed, and the model's cache is cleared
between calls, so repeated runs give identical output.

\subsection{The Pipeline}
\label{subsec:reasoning-pipeline}

The sequence is fixed. The clinical state is built; escalation is decided; retrieval runs;
the prompt is assembled; the model generates an answer; the answer is parsed into a
structured form; and the answer is checked against the state. If the check finds a serious
fault, the answer is sent back once with a single correction request. Sending one complaint
at a time was found to work better than sending several together.

The output is a ranked list of diagnoses. Each carries a likelihood, the evidence for and
against it, and the limits of that evidence. Alongside it the agent reports what
information is missing, how confident it is, and a recommended next step. The system
recommends investigations; it does not prescribe treatment.

\subsection{Checking the Answer}
\label{subsec:reasoning-safety}

The model's answer is checked against the clinical state by a set of rules. Faults are
sorted into two levels. A \emph{serious} fault --- citing a test that was never performed,
inventing a finding, misreading a value, or using an unrelated negative as an argument
against a diagnosis --- causes the differential to be withheld entirely. A \emph{minor}
fault, such as untidy wording, is delivered with a warning attached.

Withholding rather than displaying with a caveat is a deliberate choice. A differential
that cites tests which were never run is not a slightly flawed answer; it is an answer built
on evidence that does not exist, and showing it beside a warning invites the reader to use
it anyway.

These rules are exercised by 158 automated tests, grouped by the property each one
protects. Tab.~\ref{tab:safetytests} lists the largest groups. The tests use scripted
responses and need no model, so they run in seconds and are independent of which model is
installed.

\begin{table}[H]
  \centering
  \small
  \begin{tabular}{lr}
    \toprule
    \textbf{Property} & \textbf{Tests} \\
    \midrule
    Hallucination rejection      & 21 \\
    Retrieval grounding          & 19 \\
    Absent is not normal         & 16 \\
    Malformed output rejection   & 15 \\
    Escalation policy            & 11 \\
    Conflict detection           & 9 \\
    Confidence calibration       & 9 \\
    Other properties (11 groups) & 58 \\
    \midrule
    \textbf{Total}               & \textbf{158} \\
    \bottomrule
  \end{tabular}
  \caption{Safety tests by property. All pass.}
  \label{tab:safetytests}
\end{table}

\subsection{Does Retrieval Help?}
\label{subsec:reasoning-ab}

Retrieval was added after the rest of the layer was working, and the version without it was
frozen so the two could be compared fairly. Five cases were run twice: once without
retrieval and once with it. Everything else was held constant.

Three checks confirm the comparison is sound. The no-retrieval run reproduced the frozen
version exactly on all five cases; the escalation decisions were identical in both runs, as
they must be since they are computed before the model; and repeated runs gave identical
output. Tab.~\ref{tab:ab} gives the result.

\begin{table}[H]
  \centering
  \small
  \begin{tabular}{llrrl}
    \toprule
    \textbf{Case} & \textbf{Passages found} & \textbf{Errors without} &
    \textbf{Errors with} & \textbf{Effect} \\
    \midrule
    missing      & 2 & 0 & 1 & worse \\
    conflict     & 1 & 1 & 1 & error replaced \\
    concordant   & 0 & 0 & 0 & no retrieval \\
    reassuring   & 0 & 1 & 1 & no retrieval \\
    not assessed & 3 & 4 & 2 & better \\
    \midrule
    \textbf{Total} & & \textbf{6} & \textbf{5} & \\
    \bottomrule
  \end{tabular}
  \caption{Validator errors with and without retrieval, on five designed cases.}
  \label{tab:ab}
\end{table}

The totals are almost unchanged, and the aggregate hides what actually happened. Retrieval
fixed two faults: a normal troponin that had been described as elevated, and an unexamined
organ that had been left out of the list of missing information. It also created a new
one.

\begin{keystatement}
Giving a small model reference text creates a new way for it to fabricate: it attributes
the reference's general statements to the specific patient.
\end{keystatement}

In two of the three cases where retrieval returned anything, the model took a sentence from
a retrieved passage and listed it as evidence about this patient. In one case it cited
``lung sliding (low specificity)'', a phrase from the pneumothorax unit. In another it
cited ``absence of sonographic signs of RV overload or dysfunction'', taken word for word
from the cardiac unit. Neither was an observation about the patient in front of it.

Every instance was caught by the checking rules and the differential was withheld. This is
the clearest evidence in the project for keeping the safety layer independent of the model:
retrieval opened a new way to fail, and the checks closed it without being changed.

Three limits apply to this experiment. Two of the five cases retrieved nothing, so they say
nothing about retrieval; the case with an entirely normal patient produces a query too thin
to match anything, which is a gap in how queries are built. The improvement on the third
case is partly an artefact, because the model changed its diagnosis rather than reasoning
better about the same one. And the corpus unit describing what a negative FAST misses does
not itself surface for that question, because TF--IDF matches words rather than meaning.

The conclusion is narrow and worth stating plainly: on these five cases, retrieval did not
make the system better. It changed which errors occurred.

\subsection{Limitations}
\label{subsec:reasoning-limits}

The safety layer behaved as specified on five designed cases and on 158 automated tests.
That is the whole of what has been shown. The reasoning quality of the local model remained
limited, and the system was not evaluated against clinical ground truth: there is no
reference standard here saying what the correct differential was for any case. The five
cases were written to exercise particular failure modes, not sampled from a population, so
no rate of anything can be inferred from them.

This is decision support and not a diagnostic device. It has not been validated against
patient outcomes and is not suitable for clinical use.
